The design is a structure that can be reasoned about through the meaning that it establishes. In this context, reasoning means to ask meaningful questions about it to extract a meaningful answer. As such, if the design is established in a structured and unambiguous form, an engine can automatically reason about the world.

\subsection{Meaning of Failures}

Reasoning about the world is a familiar process to developers because it has always been a part of software development. This is because software systems have always been the design of a closed semantic world. The process of reasoning about the world is asking meaningful questions about it to get meaningful answers.

When a meaningful question cannot be answered, it exposes missing meaning and in a closed semantic world, this is the meaning of a failure. As such, in this document, the world failure means an inability by the engine to answer a meaningful question about the world.

\begin{figure}[htbp]
    \centering
    \includesvg[
        width=0.85\linewidth
    ]{assets/reasoning_loop.svg}
    \caption{Automating Failure Diagnosis by Reasoning about World to Answer Meaningful Questions}
    \label{fig:meaning_of_failures}
\end{figure} 

The word failure has historical baggage because it has often been used to reference mechanical failure experienced by the implementation of a design on a substrate (which is one possible result of semantic failure). So I want to make that distinction and clarify that the use of the word failure in this document is not to be confused with mechanical failure. It is worth highlighting this explicitly because understanding what failures are in a semantic world is the heart of this solution and the way it frames software systems.

In earlier sections, I established the conditions under which semantic invalidity exists. The persistence of the semantic invalidity can be captured through either an inability to answer a meaningful question about the world or an inability to identify the meaningful question entirely. In the following sections, I will describe how semantic invalidity is eliminated through reasoning by expanding the semantics to identify or answer meaningful questions to eliminate failures. 

\subsection{Design Abstraction Language}

To make the reasoning possible, the design must be defined in a Design Abstraction Language (DAL). The DAL is a declarative language that is used fully establish the closed semantic world. The scope of the DAL includes all the meaning needed to fully define the world (behaviors, transitions, participants, attributes, invariants, etc). 

Later in the document, I will describe a semantic toolkit that can be used establish the meaning of a world. This meaning is automatically translated into a closed semantic world defined in the Design Abstraction Language (DAL) and the engine reasons about the world using the structure until the meaning is correct. The established meaning can then be added to the semantic toolkit and the process repeats.

\subsection{Types of Reasoning}

An engine can use the DAL definition to automatically reason about the world. When the engine is unable to answer a meaningful question about the world, the failure is automatically diagnosed because it identifies where the semantics have to expand through root cause analysis. 

Figure~\ref{fig:meaning_of_failures} demonstrates the general structure of how the engine can automatically reason and expand the meaning of the world by learning new semantics. In the following sections, I will demonstrate how this reasoning is applied to learn from the requirements, the design own meaning and reality.

\subsubsection{Reasoning about Requirements}

By answering meaningful questions about the requirements, the engine can reason about the world and identify where the meaning of the world has to expand through root cause analysis. This feedback loop converges when the meaning of the world fully satisfies the requirements.

\begin{figure}[htbp]
    \centering
    \includesvg[
        width=0.85\linewidth
    ]{assets/requirement_reasoning.svg}
    \caption{Reasoning to Verify Requirement Satisfaction}
    \label{fig:requirement_reasoning}
\end{figure}

This is possible because the requirements and the design share meaning. If the meaning of design satisfies all the requirements, then the necessary conditions to verify the internal consistency of the design will exist. It also becomes clear that if the requirements themselves are ambiguous, then the design will inherit this ambiguity.

However, this also implies that if the requirements are specified in a structured and unambiguous way, the design can be automatically generated from the requirements. When the completeness of the requirements converges on its limit, it becomes the meaning of the design itself. This will eliminate the need to reason about the requirements to establish the necessary meaning in the world because the meaning itself will be specified and used to construct the world.

In this process, building the design itself becomes an automated reasoning process, where the meaning is used to construct the world and the engine identifies the missing meaning that prevents internal consistency. Through this feedback loop, the meaning keep getting refined until the design is internally consistent.

\subsubsection{Reasoning about Internal Consistency}

As the previous section described, the closed semantic world is automatically constructed using the specified meaning. As a result, using this meaning, the engine can then reason about the internal consistency of the design. 

\begin{figure}[htbp]
    \centering
    \includesvg[
        width=0.95\linewidth
    ]{assets/reshape_design.svg}
    \caption{Automated Reasoning to Reshape Narrative}
    \label{fig:reshape_design}
\end{figure}

Through the meaning of the world, it is possible to identify semantically invalid narratives. For example, if a narrative gets a book from the basket but the world allows narrative to be exhibited when the basket is empty, then this will result in a failure as predicted by the invariant. This can be resolved by selecting a semantically valid narrative to prevent the invalid narrative from being selected. This process is demonstrated in a general way in Figure~\ref{fig:invariant_path}. 

However, if the necessary meaning doesn't exist in the world to be able to reason about the designs internal consistency, then the failure is automatically diagnosed and the subsequent root cause analysis expands the meaning of the world. For example, if the invariant specifies that the basket size has to be at least one before you try to get a book from it, and the design doesn't know what basket size means, it will be unable to reason about it. In this way, the meaning of the world will expand until the engine can eliminate all semantically invalid narratives.

While this is a simple example, it highlights a more important point that inconsistences in the constructed narratives will be identified by the reasoning. If there is no narrative that adds a book to the basket, then is the get book from basket narrative actually needed or is the behavior to add a book to the basket missing? The engine can reason about this because it recognize that the world state never modifies the basket size, so you can never get a book from the basket. 

In this sense, this process will actually expand the requirements of the design until it is complete and the internal consistency can be verified.

\subsubsection{Reasoning about Reality}

Since the design establishes its assumptions unambiguously, the engine can reason about reality to determine when the reality refutes the designs assumptions. Since it will be unable to explain why reality refuted the designs assumptions, this will result in automated failure diagnosis and the designs semantics will expand to explain why. Mechanical failures are just one way in which reality can refute the designs assumptions, any assumptions made by the design like timing requirements can also be refuted.

\begin{figure}[htbp]
    \centering
    \includesvg[
        width=0.85\linewidth
    ]{assets/ideal_learning.svg}
    \caption{Ideal learning loop by reasoning about execution using design meaning}
    \label{fig:ideal_learning}
\end{figure}

In this process, the design learns how to function effectively in its operational environment and becomes progressively more intelligent. This happens through the acquired environmental invariants which expand the deigns awareness of its environment.

However, in order to practically enable this process of learning from reality, the execution of the design must be lossessly preserved so that the engine can reason about it using the model's assumptions. In the next section, I will discuss the practical issues that must be considered.

\subsection{Practical Considerations for Reasoning about Reality}

In order for the engine to reason about the execution to learn new semantics, the following conditions must be met:
\begin{itemize}
    \item The implementation and its execution must represent the design semantics faithfully.
    \item The execution must be losslessly preserved for the reasoning to be deterministic.
\end{itemize}

If the semantic gap between the implementation and the design is eliminated and the execution is losslessly preserved, it would enable a deterministic learning loop. More importantly, as a result of this, any meaningful question about the world can be automatically answered or an automatic process would collect the necessary diagnostic information needed to learn how to answer the question.

In the following sections, I will explain how a Computable Semantic Model (CSM) can be used to eliminate the semantic gap between the design and the implementation. Then I will explain how domain specific compression can be used to losslessly preserve all the narratives experienced in the world so the engine can deterministically reason about it and finally I will explain how all of this works to fully automate systems management.